\documentclass[a4paper]{arrowhead}
\usepackage[T1]{fontenc}
\usepackage{longtable,booktabs,xurl}
\setlength{\emergencystretch}{3em}
\setlength{\parindent}{0pt}
\setlength{\parskip}{5pt}
\begin{document}
\ArrowheadTitle{CertificateAuthority (Rust)}
\ArrowheadType{HTTPS Interface Design Description}
\ArrowheadTypeShort{Rust}
\ArrowheadVersion{5.2.2}
\ArrowheadDate{2026-10-17}
\ArrowheadAuthor{Implementation working draft}
\ArrowheadStatus{DEVELOPMENT}
\ArrowheadContact{Repository maintainers}
\ArrowheadFooter{Eclipse Arrowhead -- Rust support development port}
\ArrowheadSetup
\section{CertificateAuthority}
\subsection{System overview}
Native Rust development profile for local Java source \path{JavaSupportRuntime/}. One independently configured executable, identity, port and state file per role. No Java, Go or Python process implements the service. Some essential operations are explicitly unavailable.
\subsection{Design and trust boundaries}
Tokio/Hyper/Axum with rustls enforces mutual TLS. Certificate chain, validity, local CN/profile and startup-configured optional CRLs are checked. CRLs need process restart to refresh; without CRLs no revocation enforcement is claimed. Tokens are verified with Authentication and bound to the TLS identity. Business calls require ConsumerAuthorization; management requires a verified sysop or explicit broad whitelist. CA provisioning instead applies Java-style named certificate-principal rights. Outbound authorities are pinned to expected CN/profile before credentials are sent, in addition to chain, SAN and time validation.
\subsection{Persistent state}
Independent locked atomic JSON replaces SQL. Clone/validate/fsync/rename commits prevent partially published local transactions. Rename is the commit point; directory-sync failure is an uncertain durable commit, not a rollback. Owner-only keys/password files are separate. Network effects are not distributed transactions; interrupted starts and bridge operations require reconciliation. No automatic Java state import, clustering or deployment is delivered.
\subsection{Implemented behavior}
Native CSR verification/signing, profiles, revocation, PEM CRLs, trusted keys
\subsection{Open release gates}
Explicit SANs only (no CSR extension copying); CRL file publication/reload, OCSP/HSM/HA, full validity/DTO normalization. All systems remain development profiles, not certified interchangeable Java replacements. See \path{COMPATIBILITY.md} for detailed DTO fields and restrictions; the Java sources remain the authoritative contract.
\subsection{Produced and consumed services}
Advertised implemented HTTPS definitions: \path{certificateAuthorityManagement}, \allowbreak \path{certificateSigning}, \allowbreak \path{certificateValidation}, \allowbreak \path{echo}, \allowbreak \path{generalManagement}.\par
Authorities consumed as applicable: Authentication, ServiceRegistry, ConsumerAuthorization and explicitly configured pinned adapter/provider origins. Only implemented operations are registered. Service versions are 1.0.0 except Gatekeeper/GateTunnel 5.2.2. No AMQP sysMonitor.observe implementation is advertised.
\subsection{HTTPS operations}
JSON unless raw configuration bytes, PEM CRL or echo text. Bodies: 1 MiB; remote responses: 8 MiB. Header: \texttt{Bearer IDENTITY-TOKEN//token}. Unknown routes return 404, absent dependencies 503. Unported operations return 501 after applicable access checks and are excluded from registration. Route presence is not semantic parity.
\begin{longtable}{p{0.09\textwidth}p{0.59\textwidth}p{0.20\textwidth}}\toprule
Method & Path & Access / status\\\midrule\endhead
\texttt{GET} & \path{/certificate-authority/general/mgmt/get-config} & management\\
\texttt{POST} & \path{/certificate-authority/general/mgmt/logs} & management\\
\texttt{GET} & \path{/certificate-authority/monitor/echo} & certificate\\
\texttt{POST} & \path{/certificate-authority/sign} & certificate\\
\texttt{GET} & \path{/certificate-authority/crl} & certificate\\
\texttt{POST} & \path{/certificate-authority/checkCertificate} & certificate\\
\texttt{POST} & \path{/certificate-authority/checkTrustedKey} & certificate\\
\texttt{GET} & \path{/certificate-authority/mgmt/certificates} & certificate\\
\texttt{DELETE} & \path{/certificate-authority/mgmt/certificates/{id}} & certificate\\
\texttt{POST} & \path{/certificate-authority/mgmt/revocations} & certificate\\
\texttt{GET} & \path{/certificate-authority/mgmt/keys} & certificate\\
\texttt{PUT} & \path{/certificate-authority/mgmt/keys} & certificate\\
\texttt{DELETE} & \path{/certificate-authority/mgmt/keys/{id}} & certificate\\
\bottomrule\end{longtable}
Error envelope: errorMessage, errorCode, exceptionType, origin and timestamp. Full Java status, DTO and validation equivalence requires differential testing.
\section{References and qualification}
See \path{contracts.json} for local Java SHA256 hashes/mappings, \path{COMPATIBILITY.md} for DTO details and gaps, and \path{VERIFICATION.md} for tests actually run. No upstream release-tag provenance, complete replacement or real-cloud interoperability is asserted.\end{document}
